\section*{Hoofdstuk \ref{ch:b2:amines} --- Aminen en stikstofverbindingen}

\begin{solution}{exo:b2:amines:1}
Primair; secundair; tertiair; quaternair ammoniumzout; primair (\omterm{def:b2:huckel:aromatic}{aromatisch}).
\end{solution}

\begin{solution}{exo:b2:amines:2}
Ethaanamide $<$ aniline $<$ ammoniak $<$ methylamine ($\mathrm pK_a$ van de geconjugeerde zuren
$-0.6$, 4.60, 9.26, 10.66).
\end{solution}

\begin{solution}{exo:b2:amines:3}
\textit{N}-Methylcyclohexaanimine, \ce{C6H10=N-CH3}: amine en keton in een zwak zuur milieu
(pH 4--5), met verwijdering van het gevormde water (droogmiddel of azeotropische destillatie).
\end{solution}

\begin{solution}{exo:b2:amines:4}
Broombenzeen; benzonitril; fenol; 4-hydroxyazobenzeen \ce{C6H5-N=N-C6H4-OH} (koppeling para ten opzichte van de
groep \ce{O-} van het fenolaat).
\end{solution}

\begin{solution}{exo:b2:amines:5}
$\mathrm pK_a = 9.80$: bij pH 7 is de verhouding $[\ce{(CH3)3NH+}]/[\ce{(CH3)3N}] = 10^{2.8} \approx 630$: de
geprotoneerde vorm overheerst; citroensap (pH rond 2) protoneert praktisch alles. Een
organische oplossing van een amine schudden met verdund waterig zuur brengt het amine, als zijn ammoniumion,
over naar het water; een base maakt het weer vrij.
\end{solution}

\begin{solution}{exo:b2:amines:6}
1-(Cyclohex-1-een-1-yl)pyrrolidine. Na additie en verlies van water draagt het stikstofatoom een
positieve lading (een iminiumion) en heeft het geen waterstofatoom om af te staan; in plaats daarvan gaat een proton verloren van het naburige
koolstofatoom, waarbij de binding \ce{C=C} ontstaat.
\end{solution}

\begin{solution}{exo:b2:amines:7}
Benzaldehyde met propaan-2-amine, of propanon met benzylamine, telkens met \ce{NaBH3CN} bij pH
rond 6.
\end{solution}

\begin{solution}{exo:b2:amines:8}
Propiofenon (1-fenylpropaan-1-on), \ce{C6H5COCH2CH3}: één additie geeft het imine-anion,
dat tot het keton gehydrolyseerd wordt.
\end{solution}

\begin{solution}{exo:b2:amines:9}
Kaliumftaalimide + 1-broomhexaan: \textit{N}-hexylftaalimide; daarna maakt hydrazine (of hydrolyse)
hexaan-1-amine vrij. Het gealkyleerde ftaalimidestikstofatoom heeft geen beschikbaar vrij elektronenpaar (het is
gedelokaliseerd over twee carbonylgroepen) en geen binding \ce{N-H} meer: het kan geen tweede keer gealkyleerd worden.
\end{solution}

\begin{solution}{exo:b2:amines:10}
Broomwater bromeert aniline op de drie posities ortho en para ten opzichte van de aminogroep:
2,4,6-tribroomaniline. \omterm{def:b2:amines:diazonium}{Diazotering} en daarna \ce{H3PO2} vervangt de aminogroep door H: de drie
broomatomen staan nu 1,3,5 ten opzichte van elkaar.
\end{solution}

\begin{solution}{exo:b2:amines:11}
4-Nitroaniline + \ce{NaNO2} + \ce{HCl} bij 0--\qty{5}{\degreeCelsius}: 4-nitrobenzeendiazoniumchloride.
Koppeling met fenol in zwak basische oplossing (fenolaat, sterk geactiveerd) op de
positie para ten opzichte van \ce{O-}: 4-[(4-nitrofenyl)diazenyl]fenol, \ce{O2N-C6H4-N=N-C6H4-OH}, een oranje
kleurstof.
\end{solution}

\begin{solution}{exo:b2:amines:12}
De trimethylammoniumgroep is een slechte, omvangrijke vertrekkende groep en maakt de $\beta$-waterstofatomen aan de
kant van \ce{CH3} zuurder en beter bereikbaar; de base haalt een waterstofatoom weg van het minst
gesubstitueerde koolstofatoom: but-1-een, het minst gesubstitueerde alkeen (hofmannproduct).
\end{solution}

\begin{solution}{pb:b2:amines:1}
\textbf{1.} $\mathrm{{}^-O_3S{-}C_6H_4{-}NH_3^+}$: de zure sulfonzuurgroep heeft het amine geprotoneerd.
\textbf{2.} Het zwitterion lost slecht op in water; carbonaat deprotoneert de ammoniumgroep,
wat het oplosbare natriumsulfanilaat geeft.
\textbf{3.} \ce{NaNO2 + HCl -> HNO2 + NaCl}.
\textbf{4.}
\[
  \mathrm{{}^-O_3S{-}C_6H_4{-}NH_2 + HNO_2 + H^+ \to {}^-O_3S{-}C_6H_4{-}N_2^+ + 2\,H_2O} .
\]
\textbf{5.} Het diazoniumzout ontleedt (tot het fenol en \ce{N2}) als het warm wordt; salpeterigzuur
ontleedt ook.
\textbf{6.} $5.00/173.19 = \qty{28.9}{mmol}$: \qty{28.9}{mmol} \ce{NaNO2}, \qty{1.99}{g}.
\textbf{7.} Een overmaat salpeterigzuur zou met de koppelingspartner reageren; ze wordt aangetoond met jodide-zetmeelpapier
en vernietigd met ureum of amidosulfonzuur.
\textbf{8.} Zijn ring wordt sterk geactiveerd door de dimethylaminogroep, een krachtige donor.
\textbf{9.} Para ten opzichte van \ce{N(CH3)2}: ortho/para-sturing, en de orthoposities worden gehinderd door de
methylgroepen.
\textbf{10.} $\mathrm{{}^-O_3S{-}C_6H_4{-}N{=}N{-}C_6H_4{-}N(CH_3)_2}$.
\textbf{11.} In sterk zuur wordt de dimethylaminogroep geprotoneerd: \ce{-NH(CH3)2+} desactiveert de ring
en de koppeling stopt.
\textbf{12.} Het product vormt zich als zijn geprotoneerde, gedeeltelijk oplosbare zure vorm; in base slaat het natriumsulfonaat,
minder oplosbaar in de zoute oplossing, neer.
\textbf{13.} Zijn positieve lading is gedelokaliseerd over de ring en het eindstandige stikstofatoom is een zwak
elektrofiel centrum; alleen ringen die geactiveerd zijn door \ce{-O-}, \ce{-OH} of \ce{-NR2} reageren.
\textbf{14.} Geeloranje in base, rood in zuur.
\textbf{15.} Op een stikstofatoom van de azogroep (de geprotoneerde vorm wordt gestabiliseerd door conjugatie met
de dimethylaminogroep).
\textbf{16.} Ongeveer $\mathrm pK_a \pm 1$: van ruwweg 3.1 tot 4.4, rond $\mathrm pK_a = 3.47$.
\textbf{17.} Twee ringen en de azogroep vormen een lang geconjugeerd systeem: de kloof tussen het hoogste
bezette en het laagste lege $\pi$-niveau is klein (\cref{prop:b2:huckel:gap}) en de absorptie
ligt in het zichtbare.
\textbf{18.} Protonering verandert de ladingsverdeling in het $\pi$-systeem en verkleint de kloof
nog verder: de absorptie verschuift naar langere golflengten, de kleur van geel naar rood.
\textbf{19.} \qty{173.19}{g/mol}.
\textbf{20.} \qty{327.33}{g/mol}.
\textbf{21.} \qty{28.87}{mmol}.
\textbf{22.} $0.02887 \times 327.33 = \qty{9.45}{g}$.
\textbf{23.} Ontleding van het diazoniumzout, onvolledige koppeling, oplosbaarheid van het product in
de moederloog en het waswater.
\textbf{24.} $0.80 \times 9.45 = \boldsymbol{\qty{7.56}{g}}$ methyloranje.
\end{solution}
